% !TEX root = ../memoria.tex
\capitulo{2}{Objetivos del proyecto}

Este proyecto persigue construir una aplicación web que permita consultar, en lenguaje natural, el estado y la predicción del tráfico en distintas ciudades de España, apoyándose en modelos de redes neuronales de grafos (\textit{Graph Neural Networks}, GNN) y en el protocolo \textit{Model Context Protocol} (MCP) como capa de comunicación entre un modelo de lenguaje y los datos y modelos del sistema. A continuación se detallan los objetivos, distinguiendo entre el general, los derivados de los requisitos del software y los de carácter técnico.

\section{Objetivo general}

Desarrollar un sistema de apoyo a la movilidad urbana que responda a preguntas sobre el tráfico de ciudades españolas mediante una interfaz web, utilizando un modelo de predicción basado en GNN expuesto a través de un servidor MCP.

\section{Objetivos del software}

\begin{itemize}
  \item Permitir al usuario formular preguntas en lenguaje natural sobre el estado actual del tráfico en una ciudad y franja horaria concretas.
  \item Ofrecer predicciones de tráfico a corto plazo, indicando el horizonte temporal de la predicción.
  \item Presentar las respuestas de forma comprensible, indicando la ciudad, el momento y el nivel de congestión estimado.
  \item Soportar varias ciudades españolas, de modo que añadir una nueva no requiera modificar la lógica de la aplicación.
  \item Proteger los datos de la universidad: los datos originales permanecerán en el entorno donde se alojan y el modelo de lenguaje solo recibirá resultados agregados.
\end{itemize}

\section{Objetivos técnicos}

\begin{itemize}
  \item Estudiar el estado del arte de MCP, de las redes LSTM y de las GNN aplicadas a la predicción de tráfico.
  \item Analizar, limpiar y preparar los datos de tráfico proporcionados por la universidad, definiendo la estructura de grafo (sensores y conexiones viarias) necesaria para las GNN.
  \item Implementar un modelo de referencia basado en LSTM y un modelo basado en GNN, y compararlos mediante métricas de error de predicción como MAE y RMSE.
  \item Diseñar e implementar un servidor MCP en Python que exponga herramientas de consulta del estado del tráfico y de predicción, integrando los modelos entrenados.
  \item Desarrollar una aplicación web con un cliente MCP que conecte un modelo de lenguaje con el servidor, de forma que sea este quien decida qué herramienta invocar.
  \item Evaluar el sistema completo, tanto la calidad de las predicciones como el comportamiento de las respuestas a preguntas representativas de los usuarios.
\end{itemize}

\section{Objetivos de gestión y documentación}

\begin{itemize}
  \item Aplicar una metodología ágil con iteraciones (\textit{sprints}) de dos semanas, con objetivos revisados al final de cada una.
  \item Mantener el trabajo documentado y versionado en un repositorio de GitHub, con \textit{issues} y \textit{milestones} asociados a cada iteración.
  \item Elaborar la memoria del proyecto de forma continua durante el desarrollo, recogiendo las decisiones de diseño y los resultados.
\end{itemize}
